\subsection{Unique solvability and simple SCMs}

\begin{Def}\label{def:scm_unique_solvability}
  An SCM $M = \lt \Sin, \Send, \Srnd, \CX, P, f \rt$ is called \emph{uniquely solvable} if it has a unique solution function. This means two things: (i) it has a solution function $g : \CXJ \times \CXW \to \CXV$, i.e., a measurable function
  that satisfies
  \begin{equation}\label{eq:scm_unique_solvability}
    \forall \xW\in\CXW \forall \xJ\in\CXJ: \quad g(\xJ,\xW) = f\big(\xJ, g(\xJ,\xW), \xW\big);
  \end{equation}
  (ii) all its solution functions must equal $g$, i.e., for any solution function $\tilde{g} : \CXJ \times \CXW \to \CXV$,
%  we have that for $P$-almost all $\xW\in\CXW$, for all $\xJ\in\CXJ$:
  we have that for all $\xW\in\CXW$, for all $\xJ\in\CXJ$:
  $$g(\xJ,\xW) = \tilde{g}(\xJ,\xW).$$
  %We will call such a solution function \emph{unique} since it is unique up to $P$-null sets.
\end{Def}
The following result gives two useful properties of unique solvability. 
The first provides an equivalent formulation that makes it easier to check whether an SCM
is uniquely solvable, the second provides an important consequence regarding
the Markov kernels of the SCM.
\begin{Thm}\label{thm:scm_unique_solvability}
  Let $M = \lt \Sin, \Send, \Srnd, \CX, P, f \rt$ be an SCM.
  \begin{enumerate}
    \item $M$ is uniquely solvable if and only if 
      for all $\xW\in\CXW$, for all $\xJ \in \CXJ$, the equation
      $$\xV = f(\xJ,\xV,\xW)$$
      has a unique solution for $\xV\in\CXV$.
%    \item If $M$ is uniquely solvable, then for all $\xJ\in\CXJ$, $M_{\doit(\XJ=\xJ)}$ is uniquely solvable.
    \item If $M$ is uniquely solvable, it has a unique Markov kernel 
      $$P_{M}(\XV,\XW \mid\doit(\XJ)) = (g,\id_{\CXW})_* \lp \bigg( \bigotimes_{w\in W} P_w(X_w) \bigg) \otimes \delta(X_J|X_J) \rp $$
      where $g$ is the unique solution function of $M$.
    \item If $M$ is uniquely solvable, the distribution of a potential outcome $X^{\doit(\xJ)}_{V\cup W}$ for input $\xJ\in\XJ$ is given by $P_{M}(\XV,\XW \mid \doit(\XJ=\xJ))$.
  \end{enumerate}
\end{Thm}
\begin{proof}
  \begin{enumerate}
    \item 
    ``$\implies$'': Let $g : \CXJ \times \CXW \to \CXV$ be a solution function for $M$. 
    Let $\xi_W\in\CXW$, $\xi_J\in\CXJ$. The equation $\xV = f(\xV,\xi_J,\xi_W)$ does have a solution for
    $\xV$ (indeed, $g(\xi_J,\xi_W)$ is such a solution). Suppose its solution is not unique, i.e., it has another
    solution $\tilde{x}_v \ne g(\xi_J,\xi_W)$. Consider the modified function
      $$\tilde{g} : \CXJ \times \CXW \to \CXV : (\xJ,\xW) \mapsto \begin{cases}
        \tilde{x}_V & \xJ = \xi_J \land \xW = \xi_W \\
        g(\xJ,\xW)  & \text{otherwise}
      \end{cases}$$
    $\tilde{g}$ is measurable (because $g$ is measurable), satisfies \eref{eq:scm_unique_solvability}, and hence it provides a solution function. 
    However, $\tilde{g} \ne g$, which contradicts the assumed unique solvability.

    ``$\impliedby$'':
    This boils down to proving that the measurability of $f$ and the uniqueness of the solutions implies the measurability of the solution function.
    %The proof makes use of some more advanced results in measure theory and can be found as part of the proof of Theorem 3.6 in \cite{BFPM21}.
    We exploit that we are dealing with standard measurable spaces.
    Define the function $g : \CXJ \times \CXW \to \CXV$ by letting $g(\xJ,\xW)$ be the (unique) solution $\xV$ of the
    equation $\xV = f(\xV,\xJ,\xW)$, with $\xV\in\CXV$. The graph of this function is
    $$\mathrm{graph}(g) = \{(\xJ,\xW,\xV) \in \CXJ \times \CXW \times \CXV : \xV = g(\xJ,\xW)\}.$$
    By assumption, we have that $\xV = g(\xJ,\xW) \iff \xV = f(\xV,\xJ,\xW)$ for all $x \in \CX$. Hence,
    $$\mathrm{graph}(g) = \{(\xJ,\xW,\xV) \in \CXJ \times \CXW \times \CXV : \xV = f(\xV,\xJ,\xW)\}.$$
    Defining the function 
    $$h : \CXJ \times \CXW \times \CXV \to \CXV \times \CXV : (\xJ,\xW,\xV) \mapsto (\xV, f(\xV,\xJ,\xW))$$
    and the diagonal $\Delta = \{(\xV,\xV) : \xV \in \CXV\} \subseteq \CXV^2$, this shows that
    $$\mathrm{graph}(g) = h^{-1}(\Delta).$$
    Since $h$ is measurable and $\Delta$ is a measurable set (because $\CXV$ is Hausdorff), $h^{-1}(\Delta)$ is a measurable set.
    By \cite[14.12]{Kec95}, because all spaces are (isomorphic to) Borel spaces, the fact that $\mathrm{graph}(g)$ is a measurable set implies that $g$ is a measurable function. Hence $g$ is a solution function. The unique solvability of $M$ follows since this is the only possible solution function.
% perhaps Fremlin 418R

%\item
%  This easily follows by using the equivalent formulation of unique solvability from the first point.
%  If $M$ is uniquely solvable, then
%  for $P$-almost all $\xW\in\CXW$, for all $\xJ \in \CXJ$, the equation
%  $$\xV = f(\xJ,\xV,\xW)$$
%  has a unique solution for $\xV\in\CXV$.
%  This implies that we can swap the quantifiers, i.e., the equation has a unique solution for $\xV\in\CXV$ for all $\xJ \in \CXJ$, for $P$-almost all $\xW\in\CXW$.
%  This then shows that for all $\xJ\in\CXJ$, the intervened SCM $M_{\doit(\XJ=\xJ)}$ is uniquely solvable (again by the first point).
\item
  Assume $M$ to be uniquely solvable and let $g : \CXJ \times \CXW \to \CXV$ be its unique solution function.
  The push-forward $(g,\id_{\CXW})_*(P)$ of the exogenous distribution $P$ of $M$ (interpreted as a constant 
  Markov kernel from $\CXJ$ to $\CXW$) provides a Markov kernel for $M$.
  
  Let $K(\XV,\XW \mid \XJ)$ denote any Markov kernel for $M$ corresponding to a solution $X : \C{U} \times \CXJ \to \CX$.
  Pick any $\xJ\in\CXJ$. 
  We have to show that $K(\XV,\XW \mid \XJ=\xJ)$ is a unique distribution.
  The distribution $K(\XV,\XW \mid \XJ=\xJ)$ is that of $X_{V\cup W}^{\doit(\xJ)}$, which is a potential outcome of $M$.
  Since $X_{V\cup W}^{\doit(\xJ)}$ is a potential outcome of $M$, we have that $\XW^{\doit(\xJ)} \sim P$ and
    $$\XV^{\doit(\xJ)} = f(\xJ,\XV^{\doit(\xJ)},\XW^{\doit(\xJ)}) \quad\text{a.s.}$$
  This implies that
    $$\XV^{\doit(\xJ)} = g(\xJ,\XW^{\doit(\xJ)}) \quad\text{a.s.}$$
   By modifying the random variable $X_{V\cup W}^{\doit(\xJ)}$ on a null set, we can obtain a random variable $Y_{V\cup W}^{\doit(\xJ)}$ such that we get equality everywhere:
  $$Y_V^{\doit(\xJ)} = g(\xJ,Y_W^{\doit(\xJ)}).$$
  This shows that the distribution of $Y_{V\cup W}^{\doit(\xJ)}$, and hence that of $X_{V\cup W}^{\doit(\xJ)}$, 
  is that of the push-forward of the exogenous distribution $P$ through the unique function
  $g_{\xJ} : \CXW \to \CXV : \xW \mapsto (g(x_J,\xW),\xW)$.
  The latter push-forward distribution is unique.
  \end{enumerate}
\end{proof}
The first equivalence states that one gets the \emph{measurability}
of the solution function for free from the measurability of the causal mechanism $f$ and 
the uniqueness of the solutions of the structural equations.
Note that for the special case of no exogenous input variables ($\Sin = \emptyset$), 
the above shows that uniquely solvable SCMs induce a unique distribution.

For SCMs whose causal mechanism is linear in terms of the endogenous variables, we can
give a sufficient condition for unique solvability, and explicitly write down the form
of their solution function.
\begin{Prp}\label{prp:linear_scms}
Let $M=\lt \Sin, \Send, \Srnd, \CX, P, f \rt$ be an SCM such that all its
endogenous variables are real-valued (i.e., $\CX_v = \RN$ for $v \in V$), and each
component of the causal mechanism is an affine combination of endogenous variables
with coefficients that may depend on exogenous variables, i.e., of the form
$$
  f_v(x) = \sum_{u\in V} B_{vu}(x_J,x_W) x_u + c_v(x_J,x_W),
$$
where $B(x_J,x_W)\in\RN^{V \times V}$ is a family of matrices (or: a matrix-valued function $\CXJ\times\CXW \to \RN^{V\times V}$) and
$c(x_J,x_W) \in \RN^V$ is a family of real-valued offsets (or: a vector-valued function $\CXJ\times\CXW \to \RN^V$).

Then, for any set $L \subseteq V$,
$(M)_{\doit(V\setminus L)}$ is uniquely solvable if and only if the matrices 
$I_{L}-B_{LL}(x_J,x_W)$ are invertible for all $x_J \in \CX_J, x_W \in \CX_W$,
where $I_L \in \RN^{L \times L}$ the identity matrix and 
$B_{LL}(x_J,x_W) \in \RN^{L \times L}$ the submatrix of $B(x_J,x_W)$.
Its unique solution function is then:
  \begin{equation*}\begin{split}
    g &: \RN^{V\setminus L} \times \RN^{J\cup W} \to \RN^L \\
      &: (x_{V\setminus L},x_J,x_W) \mapsto (I_L - B_{LL}(x_J,x_W))^{-1} \big(B_{L,{V\setminus L}}(x_J,x_W) x_{V\setminus L} + c_{L}(x_J,x_W)\big).
  \end{split}\end{equation*}
\end{Prp}
%A non-zero coefficient $B_{ij}$ for $i,j\in\C{I}$ such that $i\neq j$ corresponds with a directed edge $j\to i$ in the (augmented) graph, and a coefficient $B_{ii}=1$ for $i\in\C{I}$ corresponds with a self-cycle $i\to i$ in the (augmented) graph of the SCM. A non-zero coefficient $\Gamma_{ij}$ for $i\in\C{I}$, $j\in\C{J}$ with $\Prb_{\C{E}_j}$ a non-degenerate probability distribution over $\RN$ corresponds with a directed edge $j \to i$ in the augmented graph. A non-zero entry $(\Gamma \Gamma^T)_{ij}$ for $i,j\in\C{I}$ with $i\neq j$ such that there exists a $k\in\C{J}$ for which $\Gamma_{ik},\Gamma_{jk}\neq 0$ and $\Prb_{\C{E}_k}$ a non-degenerate probability distribution over $\RN$ corresponds with a bidirected edge $i \oto j$ in the graph of the SCM.

If an SCM is uniquely solvable, this does not necessarily mean that it is still uniquely solvable
after performing some intervention. 
To avoid the complications introduced in case solutions are absent, or present but not unique,
we will henceforth make strong assumptions regarding the existence and uniqueness of solutions. 
We (mostly) restrict our attention to a subclass of SCMs that we refer to as \emph{simple SCMs},
which are SCMs that are uniquely solvable and remain so after any hard intervention:
\begin{Def}\label{def:simple_scm}
  An SCM $M = \lt \Sin, \Send, \Srnd, \CX, P, f \rt$ 
  is called \emph{simple} if for any $T \subseteq \Sin \cup \Send \cup \Srnd$, the intervened SCM
  $M_{\doit(T)}$ is uniquely solvable.
\end{Def}
Note that this includes unique solvability of $M$ itself for $T = \emptyset$.

\begin{Eg}
Consider an SCM with structural equations
  \begin{align*}
    X_1 &= W_1 \\
    X_2 &= W_2 \\
    X_3 &= X_1 X_4 + W_3 \\
    X_4 &= X_2 X_3 + W_4 
  \end{align*}
where the $X$'s are considered real-valued endogenous variables and the $W$'s exogenous variables with domains $(-1,1) \subset \RN$.
We can solve the system of structural equations for $X$ in terms of $W$:
  \begin{align*}
    X_1 &= W_1 \\
    X_2 &= W_2 \\
    X_3 &= \frac{W_3 + W_1 W_4}{1 - W_1 W_2} \\
    X_4 &= \frac{W_2 W_3 + W_4}{1 - W_2 W_1}
  \end{align*}
Similarly, we can take any subset of the structural equations and solve it for the variables appearing on the l.h.s.\ of the equations in the subset, and obtain a unique solution. For example, only solving the structural equations for $X_3$ and $X_4$, we obtain:
  \begin{align*}
    X_3 &= \frac{W_3 + W_1 W_4}{1 - W_1 X_2} \\
    X_4 &= \frac{W_2 W_3 + W_4}{1 - W_2 X_1}
  \end{align*}
where the variables $X_1$ and $X_2$ are now considered as exogenous input variables (instead of endogenous variables).
Hence, any subset of the structural equations has a unique solution for the variables appearing on the l.h.s.\ in terms
of the remaining ones on the r.h.s., which means that this SCM is simple.
\end{Eg}

Theorem~\ref{thm:scm_unique_solvability} shows that the Markov kernels in the following definition are all uniquely defined.
\begin{Def}\label{def:simple_scm_Markov_kernels}
  If SCM $M = \lt \Sin, \Send, \Srnd, \CX, P, f \rt$ is simple, and $O \subseteq V \cup W$ is considered observed,
  then $M$ induces the \emph{observational Markov kernel $P_{M}(X_O \mid \doit(X_J))$}, and for any intervention target 
  $T \subseteq O$, it induces the \emph{interventional Markov kernel}
  $$P_{M}(X_{O\setminus T} \mid \doit(X_J),\doit(X_T)) := P_{M_{\doit(T)}}(X_{O\setminus T} \mid \doit(X_{J \cup T})),$$
  as the marginal Markov kernel of the intervened SCM $M_{\doit(T)}$.
\end{Def}
The two notations for interventional Markov kernels of simple SCMs will be used interchangeably.

The fact that these SCMs are relatively simple to deal with (because we do not have to worry about non-existence or 
non-uniqueness of solutions) motivated their name.
Even though simplicity is a strong assumption, 
the class of simple SCMs is more expressive than the class of IL-CBNs since it allows to model causal cycles, to some extent.
%causal cycles, and (ii) it allows to define counterfactuals.

\subsection{Equivalences}

Equivalence relations are ubiquitous in mathematics. They capture the notion that mathematical objects can be ``equivalent'' 
from some point of view.
\begin{Def}\label{def:equivalent_scms}
  Let $Z$ be a set and $R \subseteq Z^2$ be a relation on $Z$ (i.e., a subset of ordered pairs of $Z$). $R$ is called an \emph{equivalence relation} if 
  \begin{enumerate}
    \item $R$ is reflexive: $(a,a) \in R$ for all $a \in Z$;
    \item $R$ is symmetric: $(a,b) \in R \iff (b,a) \in R$ for all $a,b \in Z$;
    \item $R$ is transitive: if $(a,b) \in R$ and $(b,c) \in R$ then $(a,c) \in R$ for all $a,b,c\in Z$.
  \end{enumerate}
\end{Def}
In this section, we will discuss several important equivalence relations between SCMs.

\begin{Def}
  Let $M = \lt \Sin, \Send, \Srnd, \CX, P, f \rt$ and
  $\tilde{M} = \lt \Sin, \Send, \Srnd, \CX, P, \tilde{f} \rt$ be two SCMs that may
  differ only in terms of their causal mechanism. We say that
  $M$ is \emph{equivalent} to $\tilde{M}$ and write $M \equiv \tilde{M}$ if for each $v \in \Send$,
  $$\forall x \in \CX: \qquad x_v = f_v(\xJ,\xV,\xW) \iff x_v = \tilde{f}_v(\xJ,\xV,\xW).$$
  In words, $M$ and $\tilde{M}$ are considered equivalent if they at most differ in terms of their
  causal mechanisms, yet each of their structural equations has the same solutions.
\end{Def}
Equivalent SCMs are indeed equivalent for many purposes:
\begin{Prp}
\begin{enumerate}
  \item Equivalence is preserved by hard interventions:\\
    if $M \equiv \tilde{M}$ and $T \subseteq V \cup W \cup J$, then $M_{\doit(T\dots)} \equiv \tilde{M}_{\doit(T\dots)}$;
    % \item Equivalence is preserved by twinning.
  \item Unique solvability is an invariant of equivalence:\\
    if $M \equiv \tilde{M}$ then $M$ is uniquely solvable if and only if $\tilde{M}$ is uniquely solvable;
  \item Simplicity is an invariant of equivalence:\\
    if $M \equiv \tilde{M}$ then $M$ is simple if and only if $\tilde{M}$ is simple;
  \item Equivalent SCMs have the same solution functions, solutions, (potential) outcomes, and Markov kernels.
\end{enumerate}
\end{Prp}
\begin{proof}
  Exercise for the reader.
\end{proof}

\begin{Eg}\label{eg:equivalent_scm}
  The SCM with the single structural equation
  $$X = -X^3 + X + W$$
  where $X$ is endogenous, and $W$ is exogenous, is equivalent to the SCM obtained by replacing the structural equation with
  $$X = \sqrt[3]{W}.$$
  Note that $X$ no longer appears on the r.h.s..
\end{Eg}
This notion of equivalence is rather strong. 
\begin{Eg}\label{eg:self_cycle}
  The SCM with endogenous variable $X \in \RN$ and exogenous random variable $W$ with co-domain $\RN$ and structural equation
  $$X = W X + c$$
  is not equivalent to the SCM obtained by replacing the structural equation with
  $$X = \begin{cases}
          \xi & W = 1 \\
          \frac{c}{1 - W} & W \ne 1
        \end{cases}$$
  no matter how we choose $\xi$, even when $P(W=1) = 0$. 
  If one does not model interventions on exogenous random variables, weaker notions of equivalence 
  can be used that replace the ``for all'' quantifier over values of exogenous random variables by the
  ``for almost all'' quantifier (see \cite{BFPM21}).
\end{Eg}

We often make use of other notions of equivalence as well.
For simplicity of exposition, we provide the definitions only for simple SCMs 
(the general definitions are provided in \cite{BFPM21} for SCMs without exogenous input variables).
\begin{Def}\label{def:scm_obs_interv_equiv}
  Let $M = \lt \Sin, \Send, \Srnd, \CX, P, f \rt$ and
  $\tilde{M} = \ltuple \tilde{\Sin}, \tilde{\Send}, \tilde{\Srnd}, \tilde{\CX}, \tilde{P}, \tilde{f} \rtuple$ be two simple SCMs
  and $O \subseteq (\Send \cup \Srnd) \cap (\tilde{\Send} \cup \tilde{\Srnd})$ a subset.
  We say that:
  \begin{enumerate}
    \item $M$ and $\tilde{M}$ are \emph{observationally equivalent w.r.t.\ $O$} if $\CX_O = \tilde{\CX}_O$, $\CX_{J \cap \tilde{J}} = \tilde{\CX}_{J \cap \tilde{J}}$ and their marginal Markov kernels coincide: 
      $$P_{M}(X_O \mid \doit(X_J)) = P_{\tilde{M}}(X_O \mid \doit(X_{\tilde J})).$$
    This has to be interpreted as both Markov kernels being a version of a Markov kernel $\CX_{J \cap \tilde{J}} \dshto \CX_O$,
    i.e., $P_{M}(X_O \mid \doit(X_J))$ must be essentially constant in $x_{J \setminus \tilde{J}}$ and
    $P_{\tilde{M}}(X_O \mid \doit(X_{\tilde{J}}))$ must be essentially constant in $x_{\tilde{J} \setminus J}$.\footnote{In other words,
      $$P_{M}(X_O \mid \doit(X_J)) = P_{M}(X_O \mid \doit(X_{J \cap \tilde{J}})) =
      P_{\tilde{M}}(X_O \mid \doit(X_{J \cap \tilde{J}})) = P_{\tilde{M}}(X_O \mid \doit(X_{\tilde{J}})).$$}
    \item $M$ and $\tilde{M}$ are \emph{interventionally equivalent w.r.t.\ $O$} if for every subset $T \subseteq O$ the intervened SCMs $M_{\doit(T)}$ and $\tilde{M}_{\doit(T)}$ are observationally equivalent w.r.t.\ $O \setminus T$;
  \end{enumerate}
%  In case the variables in $O$ (and $J \cup \tilde{J}$) are considered observed, whereas the variables in 
%  $(V \cup W) \setminus O$ and $(\tilde{V} \cup \tilde{W}) \setminus O$ are considered latent, 
%  we sometimes omit the ``w.r.t.\ $O$'' in this terminology (i.e., we call $M$ and $\tilde{M}$ observationally equivalent, interventionally equivalent, respectively).
\end{Def}
More generally, one could define interventional equivalence not only with respect to an observed set of variables, but also with respect to a given set of interventions.

One can show the following properties of these equivalences:
\begin{Prp}\label{prp:pearls_ladder_levels_1_and_2}
  For simple SCMs $M, \tilde{M}$ and a subset $O \subseteq (\Send \cap \tilde{\Send}) \cup (\Srnd \cap \tilde{\Srnd})$:
  \begin{enumerate}
    \item If $M \equiv \tilde{M}$ then $M$ and $\tilde{M}$ are interventionally equivalent w.r.t.\ $O$.
    \item If $M$ and $\tilde{M}$ are interventionally equivalent w.r.t.\ $O$ then 
      $M$ and $\tilde{M}$ are observationally equivalent w.r.t.\ $O$.
  \end{enumerate}
\end{Prp}
\begin{comment}\begin{proof}
  Let $O'$ be the copy of $O \cap (V \cap \tilde{V})$ in $V' \cap \tilde{V}'$.
  \begin{enumerate}
    \item 
     Suppose $M \equiv \tilde{M}$. Then $M^\twin \equiv \tilde{M}^\twin$,
     and for every $T \subseteq O \cup O'$, $(M^\twin)_{\doit(T)} \equiv (\tilde{M}^\twin)_{\doit(T)}$.
     Since equivalent SCMs have the same Markov kernels, $(M^\twin)_{\doit(T)}$ and $(\tilde{M}^\twin)_{\doit(T)}$
     are observationally equivalent w.r.t.\ $(O \cup O') \setminus T$ for every $T \subseteq O \cup O'$.

  \item Suppose $M$ and $\tilde{M}$ are counterfactually equivalent w.r.t.\ $O$. 
    Then $M^\twin$ and $\tilde{M}^\twin$ are interventionally equivalent w.r.t.\ $O \cup O'$. 
  For every $T = T \subseteq O \cup O'$,
  $(M^\twin)_{\doit(T)}$ and $(\tilde{M}^\twin)_{\doit(T)}$ are observationally equivalent w.r.t. $(O \cup O') \setminus T$.

  We have to show that $M$ and $\tilde{M}$ are interventionally equivalent w.r.t.\ $O$.
  That is, we have to show that for every $S \subseteq O$, $M_{\doit(S)}$ and $\tilde{M}_{\doit(S)}$
  are observationally equivalent w.r.t.\ $O \setminus S$.

%  We first practice with a special case.
%  Take $T_2 = T_1'$ with $T_1 \subseteq O \cap (V \cap \tilde{V})$ and $T_3 = \emptyset$; then
%  $(M_{\doit(T_1)})^\twin = (M^\twin)_{\doit(T_1,T_1')}$ and $(\tilde{M}_{\doit(T_1)})^\twin = (\tilde{M}^\twin)_{\doit(T_1,T_1')}$ are observationally equivalent w.r.t. 
%  $$(O \cup O') \setminus (T_1 \cup T_2) = (O \cap (V \cap \tilde{V}) \cup O \cap (W \cap \tilde{W}) \cup (O \cap (V \cap \tilde{V}))') \setminus (T_1 \cup T_1')$$
%  $$ = O \setminus T_1 \cup ((O \cap (V \cap \tilde{V})) \setminus T_1)'$$
%  and hence w.r.t.
%  $$O \setminus T_1.$$
%  Now by Lemma~\ref{lem:obs_equiv_twin_scm}, $(M_{\doit(T_1)})^\twin$ and $M_{\doit(T_1)}$ are observationally equivalent
%  w.r.t.\ $V \setminus T_1 \cup W$ and hence w.r.t.\ $O \setminus T_1$.
%  Similarly, $(\tilde{M}_{\doit(T_1)})^\twin$ and $\tilde{M}_{\doit(T_1)}$ are observationally equivalent w.r.t\ $O \setminus T_1$.
%  Hence, by transitivity, $M_{\doit(T_1)}$ and $\tilde{M}_{\doit(T_1)}$ are observationally equivalent w.r.t\ $O \setminus T_1$.

%  Let us now consider the general case.
  Now take $S \subseteq O$ and then let $T_1 = S \cap V \cap \tilde{V}$, $T_2 = T_1'$ and $T_3 = S \cap W \cap \tilde{W}$; then $S = T_1 \cup T_3$.
  Then by assumption,
  $$((M_{\doit(T_1)})^\twin)_{\doit(T_3)} = ((M^\twin)_{\doit(T_1,T_1')})_{\doit(T_3)} = (M^\twin)_{\doit(T_1,T_2,T_3)}$$
  and 
  $$((\tilde{M}_{\doit(T_1)})^\twin)_{\doit(T_3)} = ((\tilde{M}^\twin)_{\doit(T_1,T_1')})_{\doit(T_3)} = (\tilde{M}^\twin)_{\doit(T_1,T_2,T_3)}$$
  are observationally equivalent w.r.t.\ $(O \cup O') \setminus T = (O \setminus (T_1 \cup T_3)) \cup (O' \setminus T_2)$, and hence w.r.t.\ $O \setminus S = O \setminus (T_1 \cup T_3)$.
  By Lemma~\ref{lem:obs_equiv_twin_scm}, $(M_{\doit(T_1)})^\twin$ and $M_{\doit(T_1)}$ are interventionally equivalent
  w.r.t.\ $V \setminus T_1 \cup W$. Hence
  $((M_{\doit(T_1)})^\twin)_{\doit(T_3)}$ and $(M_{\doit(T_1)})_{\doit(T_3)} = M_{\doit(S)}$ are observationally equivalent 
  w.r.t.\ $V \setminus T_1 \cup W \setminus T_3$, and hence also w.r.t.\ $O \setminus S$.
  Similarly, $((\tilde{M}_{\doit(T_1)})^\twin)_{\doit(T_3)}$ and $(\tilde{M}_{\doit(T_1)})_{\doit(T_3)} = \tilde{M}_{\doit(S)}$ are observationally equivalent w.r.t\ $O \setminus S$.
  Hence, by transitivity, $M_{\doit(S)}$ and $\tilde{M}_{\doit(S)}$ are observationally equivalent w.r.t\ $O \setminus S$.
  \item Trivial (take $T = \emptyset$).
  \end{enumerate}
\end{proof}
\end{comment}
However, the reverse implications do not hold in general. This expresses that causal modeling is more refined than probabilistic modeling.
\begin{Eg}[Observational equivalence does not imply interventional equivalence]
Consider the SCM $M$ with
$$N \sim \C{N}(\mu,\sigma^2)$$
$$C = \alpha + \beta N$$
and the SCM $\tilde{M}$ with
$$C \sim \C{N}(\tilde{\mu},\tilde{\sigma}^2)$$
$$N = \tilde{\alpha} + \tilde{\beta} C$$
These SCMs are simple and their observational distributions are respectively 
  $$P_M(N,C) = \C{N}\left(\begin{pmatrix} \mu \\ \alpha + \beta \mu \end{pmatrix}, \begin{pmatrix} \sigma^2 & \beta\sigma^2 \\ \beta\sigma^2 & \beta^2\sigma^2\end{pmatrix}\right)$$
and
  $$P_{\tilde{M}}(N,C) = \C{N}\left(\begin{pmatrix} \tilde{\alpha} + \tilde{\beta} \tilde{\mu} \\ \tilde{\mu} \end{pmatrix}, \begin{pmatrix} \tilde{\beta}^2\tilde{\sigma}^2 & \tilde{\beta}\tilde{\sigma}^2 \\ \tilde{\beta}\tilde{\sigma}^2 & \tilde{\sigma}^2\end{pmatrix}\right)$$
For certain parameter choices, they are observationally equivalent.
%If $\mu = \tilde{\alpha}\tilde{\mu}$ and $\tilde{\mu} = \alpha \mu$ and
%  $\sigma^2 = \tilde{\beta}^2\tilde{\sigma}^2$ and $\beta \sigma^2 = \tilde{\beta}\tilde{\sigma}^2$ and $\beta^2 = \tilde{\sigma}^2$ then they are observationally equivalent.
However, they are not interventionally equivalent except for very special parameter choices.
\end{Eg}

\subsection{Marginalizations}

When modeling a system, we sometimes want to ``hide'' details of a subsystem. The following operation
on SCMs that we call ``marginalization'' is a causal analogue of the marginalization of probability distributions.
The computer program analogy of the marginalization operation is to hide details within a subroutine.
Intuitively, a marginalization of an SCM over a subset of endogenous variables $L$ is obtained by first 
solving a subsystem (the structural equations
corresponding to the endogenous variables in $L$) followed by substituting the solution function of the 
subsystem into the remaining structural equations
(corresponding to the endogenous variables in $V\setminus L$).
\begin{Def}
  Let $M = \lt \Sin, \Send, \Srnd, \CX, P, f \rt$ be an SCM and $L \subseteq \Send$ such that
  $M_{\doit(V\setminus L)}$ is uniquely solvable.
  Write $K = V \setminus L$ and let $\tilde{g}_L : \CX_{\Sin \cup K} \times \CXW \to \CX_L$
  be the unique solution function for $M_{\doit(K)}$.
  Then we call $M_{\setminus L} = \ltuple \Sin, K, \Srnd, \CX_{\Sin} \times \CX_{K} \times \CX_{\Srnd}, P, \tilde{f} \rtuple$ with
  $$\tilde{f}(\xJ,x_{K},\xW) = f_{K}(\xJ,x_{K},\tilde{g}_L(\xJ,x_{K},\xW),\xW)$$
  the \emph{marginalization of $M$ over $L$}.
\end{Def}
For simple SCMs, marginalizations are obviously defined over any subset $L \subseteq \Send$.
%Note that all marginalizations of $M$ are equivalent, i.e., the marginalized causal mechanism is essentially unique.

Marginalization preserves unique solvability, as the following lemma shows.
\begin{Lem}\label{lem:marginalization_preserves_unique_solvability}
  Let $M = \lt \Sin, \Send, \Srnd, \CX, P, f \rt$ be an SCM and $L \subseteq \Send$ such that
  $M_{\doit(V \setminus L)}$ is uniquely solvable.
  If $M$ is uniquely solvable then its marginalization $M_{\setminus L}$ is uniquely solvable, and the unique solution function for
  $M_{\setminus L}$ is just $g_K = \pr_K \circ g$, where $g : \CX_J \times \CX_W \to \CX_V$ is the 
  unique solution function for $M$, $K = V \setminus L$ and $\pr_K : \CXV \to \CX_K : x \mapsto x_K$ is the canonical projection on $K$.
\end{Lem}
\begin{proof}
  Let $\tilde{g}_L : \CX_{\Sin \cup K} \times \CXW \to \CX_L$
  be the unique solution function for $M_{\doit(K)}$.
  From the properties of the solution functions, we derive that
%  $M$ is uniquely solvable means there exists a measurable function $g : \CX_{J \cup W} \to \CX_V$ such  that
  for all $x \in \CX$,
  \begin{equation*}\begin{split}
  & \begin{cases}
    x_L = g_L(x_J,x_W) \\
    x_K = g_K(x_J,x_W) \\
  \end{cases}
  \iff x_V = g(x_J,x_W) \iff x_V = f(x) 
  \iff \begin{cases}
    x_L = f_L(x) \\
    x_K = f_K(x) \\
  \end{cases} \\
  & \iff \begin{cases}
    x_L = \tilde{g}_L(x_J,x_K,x_W) \\
    x_K = f_K(x_J,x_K,x_L,x_W) \\
  \end{cases}
  \iff \begin{cases}
    x_L = \tilde{g}_L(x_J,x_K,x_W) \\
    x_K = f_K(x_J,x_K,\tilde{g}_L(x_J,x_K,x_W),x_W) \\
  \end{cases} \\
  & \iff \begin{cases}
    x_L = \tilde{g}_L(x_J,x_K,x_W) \\
    x_K = \tilde{f}(x_J,x_K,x_W) \\
  \end{cases}
  \end{split}\end{equation*}
where $\tilde{f}$ is the causal mechanism of $M_{\setminus L}$.
Hence for all $x_J \in \CX_J, x_W \in \CX_W, x_K \in \CX_K$:
$$x_K = g_K(x_J,x_W) \iff x_K = \tilde{f}(x_J,x_K,x_W).$$
Therefore, $g_K = \pr_K \circ g$ is the unique solution function for $M_{\setminus L}$, and the marginalized SCM $M_{\setminus L}$ is uniquely solvable.
\end{proof}
\begin{Rem}
This also directly implies that if $M$ is uniquely solvable and its marginalization $M_{\setminus L}$ over $L \subseteq V$ is defined, the Markov kernel of the marginalization is obtained by marginalizing the original Markov kernel:
$$P_{M_{\setminus L}}(X_{V\setminus L},X_W \mid \doit(X_J)) = P_{M}(X_{V\setminus L},X_W \mid \doit(X_J)).$$
\end{Rem}

Under certain conditions, hard interventions and marginalization commute (i.e., it does not matter in which
order we apply them).
\begin{Prp}\label{prp:interventions_marginalizations_commute}
Let $M = \lt \Sin, \Send, \Srnd, \CX, P, f \rt$ be an SCM.
    For $L \subseteq V$ such that $M_{\doit(V\setminus L)}$ is uniquely solvable,
    and a hard intervention $\doit(T\dots)$ with target $T \subseteq \Sin \cup \Srnd \cup \Send$ (of any of the three variants) such that $L \cap T = \emptyset$:
    $$(M_{\doit(T\dots)})_{\setminus L} = (M_{\setminus L})_{\doit(T\dots)}.$$
\end{Prp}
\begin{proof}
  This follows by writing out the definitions and checking commutativity of
  the operations performed on the various components of the SCM tuple one-by-one. 
\end{proof}

We will show that for simple SCMs, the marginalization operation preserves the causal semantics on the remaining variables. 
A key step is the following proposition, which gives conditions under which it does not matter whether we marginalize at once or in steps.
\begin{Prp}\label{prp:marginalization_commutes}
  Let $M = \lt \Sin, \Send, \Srnd, \CX, P, f \rt$ be an SCM and $L_1, L_2 \subseteq \Send$ such that $L_1 \cap L_2 = \emptyset$.
  If $M_{\doit(V\setminus L_1)}$ is uniquely solvable, and
  $(M_{\setminus L_1})_{\doit((V\setminus L_1)\setminus L_2)}$ is uniquely solvable, then $M_{\doit(V\setminus(L_1\cup L_2))}$ is uniquely solvable, and in that case it does not matter if we first marginalize over $L_1$ and then $L_2$, or both at once, i.e:
  $$(M_{\setminus L_1})_{\setminus L_2} = M_{\setminus (L_1\cup L_2)}.$$
\end{Prp}
\begin{proof}
  Write $K_1 = V \setminus L_1$.
  Let $\tilde{g} : \CX_{J\cup K_1} \times \CXW \to \CX_{L_1}$ be the unique solution function of $M_{\doit(K_1)}$,
  i.e., for all $x \in \CX$:
  $$x_{L_1} = \tilde{g}(x_J,x_{K_1},x_W) \iff x_{L_1} = f_{L_1}(x).$$
  Let $\tilde{f} : \CX_J \times \CX_{K_1} \times \CXW \to \CX_{K_1}$ with
  $$\tilde{f}(x_J,x_{K_1},x_W) = f_{K_1}(x_J,x_{K_1},\tilde{g}(x_J,x_{K_1},x_W),x_W)$$
  be the causal mechanism of the marginal SCM $M_{\setminus L_1}$.

  If $(M_{\setminus L_1})_{\doit(V\setminus(L_1\cup L_2))}$ is uniquely solvable, it has a unique solution function
  $\dbtilde{g} : \CXJ \times \CX_{V\setminus (L_1\cup L_2)} \times \CXW \to \CX_{L_2}$, i.e., for all $x \in \CX$:
  $$x_{L_2} = \dbtilde{g}(x_J,x_{V\setminus(L_1\cup L_2)},x_W) \iff x_{L_2} = \tilde{f}_{L_2}(x_J,x_{L_2},x_{K_1\setminus L_2},x_W)$$
  Define the function $h : \CXJ \times \CX_{V\setminus(L_1\cup L_2)} \times \CXW \to \CX_{L_1 \cup L_2}$ by
  \begin{align*}
    h_{L_1}(x_J,x_{V\setminus(L_1\cup L_2)},x_W) &= \tilde{g}(x_J,x_{K_1\setminus L_2},\dbtilde{g}(x_J,x_{V\setminus(L_1 \cup L_2)},x_W),x_W) \\
    h_{L_2}(x_J,x_{V\setminus(L_1\cup L_2)},x_W) &= \dbtilde{g}(x_J,x_{V\setminus(L_1 \cup L_2)},x_W)
  \end{align*}
  Then for all $x \in \CX$:
  \begin{equation*}\begin{split}
%    & x_{L_1\cup L_2} = f_{L_1 \cup L_2}(x) \\
%    \iff &
    & \begin{cases}
      x_{L_1} &= f_{L_1}(x) \\
      x_{L_2} &= f_{L_2}(x) 
    \end{cases} \\
    \iff &
    \begin{cases}
      x_{L_1} &= \tilde{g}(x_J,x_{K_1},x_W) \\
      x_{L_2} &= f_{L_2}(x_J,x_{K_1},x_{L_1},x_W) 
    \end{cases} \\
    \iff &
    \begin{cases}
      x_{L_1} &= \tilde{g}(x_J,x_{K_1},x_W) \\
      x_{L_2} &= f_{L_2}(x_J,x_{K_1},\tilde{g}(x_J,x_{K_1},x_W),x_W) 
    \end{cases} \\
    \iff &
    \begin{cases}
      x_{L_1} &= \tilde{g}(x_J,x_{K_1\setminus L_2},x_{L_2},x_W) \\
      x_{L_2} &= \dbtilde{g}(x_J,x_{V\setminus(L_1 \cup L_2)},x_W)
    \end{cases} \\
    \iff &
    \begin{cases}
      x_{L_1} &= \tilde{g}(x_J,x_{K_1\setminus L_2},\dbtilde{g}(x_J,x_{V\setminus(L_1 \cup L_2)},x_W),x_W) \\
      x_{L_2} &= \dbtilde{g}(x_J,x_{V\setminus(L_1 \cup L_2)},x_W)
    \end{cases} \\
    \iff &
      x_{L_1\cup L_2} = h(x_J,x_{V\setminus(L_1\cup L_2)},x_W) 
  \end{split}\end{equation*}
  where in the first equivalence we used the unique solvability of $M_{\doit(K_1)}$, in the second
  equivalence we used substitution, in the third equivalence we used the unique solvability of
  $(M_{\setminus L_1})_{\doit(V\setminus (L_1\cup L_2))}$, in the fourth equivalence we used
  substitution again, and in the fifth equivalence we used the definition of $h$.
  Therefore, $h$ is the unique solution function for $M_{\doit(V\setminus(L_1\cup L_2)}$, which must therefore
  be uniquely solvable. By checking the definition, one concludes that $(M_{\setminus L_1})_{\setminus L_2} = M_{\setminus(L_1 \cup L_2)}$.
\end{proof}
\begin{Prp}\label{prp:marginalization_preserves_simplicity}
  Let $M = \lt \Sin, \Send, \Srnd, \CX, P, f \rt$ be a simple SCM. For any $L \subseteq \Send$, 
  its marginalization $M_{\setminus L}$ is also simple.
\end{Prp}
\begin{proof}
  Let $T \subseteq J \cup W \cup V \setminus L$. 
  By Proposition~\ref{prp:interventions_marginalizations_commute}, 
  the marginalization commutes with the hard intervention:
  $$(M_{\setminus L})_{\doit(T)} = (M_{\doit(T)})_{\setminus L}.$$
  Because $M$ is simple, also $M_{\doit(T)}$ is simple (by Proposition~\ref{prp:hard_interventions_commute}), 
  and in particular it is
  uniquely solvable. From Lemma~\ref{lem:marginalization_preserves_unique_solvability} it follows that also
  its marginalization $(M_{\doit(T)})_{\setminus L}$ is uniquely solvable. 
  This means that $(M_{\setminus L})_{\doit(T)}$ is uniquely solvable. Since this
  holds for any $T \subseteq J \cup W \cup V \setminus L$, $M_{\setminus L}$ is simple.
\end{proof}
\begin{Cor}
  Let $M = \lt \Sin, \Send, \Srnd, \CX, P, f \rt$ be a simple SCM.
  For $L_1, L_2 \subseteq \Send$ with $L_1 \cap L_2 = \emptyset$,
    $$(M_{\setminus L_1})_{\setminus L_2} = (M_{\setminus L_2})_{\setminus L_1} = M_{\setminus (L_1 \cup L_2)}.$$
\end{Cor}

These commutation relations and compatibilities now allow us to give a straightforward proof that
the causal semantics are preserved under marginalization.
While this holds generally \cite{BFPM21}, we will here only prove this for simple SCMs.
\begin{Thm}\label{Thm:various_equiv_marginalization}
  Let $M = \lt \Sin, \Send, \Srnd, \CX, P, f \rt$ be a simple SCM, $L \subseteq \Send$, and
  $M_{\setminus L}$ its marginalization over $L$. Then $M$ and $M_{\setminus L}$ are observationally
  and interventionally equivalent w.r.t.\ $(V \cup W) \setminus L$.
  % counterfactual equivalence moved to a corollary later on
\end{Thm}
\begin{proof}
  %The statement that $M$ and $M_O$ are observationally equivalent w.r.t.\ $O$ is Lemma~\ref{lem:obs_equiv_marginalization}.
  Write $K = V \setminus L$.
  We first show that the marginal Markov kernels $P_{M}(X_K,X_W \mid \doit(X_J))$ and $P_{M_{\setminus L}}(X_K,X_W \mid \doit(X_J))$ are the same.
%  By Theorem~\ref{thm:scm_unique_solvability}, the observational Markov kernels $P_{M}(X_V, X_W \mid X_J)$ and $P_{M_K}(X_K, X_W \mid X_J)$ exist and are unique. Hence, the marginal observational Markov kernels $P_{M}(X_K,X_W \mid X_J)$ and $P_{M_K}(X_K,X_W \mid X_J)$ exist and are unique. 
  The former is obtained as:
  $$P_{M}(X_K,X_W \mid \doit(X_J)) = (\pr_{K\cup W} \circ (g,\id_{\CXW}))_*(P) = (g_{K},\id_{\CXW})_*(P),$$
  where $g : \CXJ \times \CXW \to \CXV$ is the unique solution function of $M$ and
  $\pr_{K\cup W} : \CX_{V \cup W} \to \CX_{K\cup W}$ is the canonical projection on the $K\cup W$ components. 
  The latter is obtained as:
  $$P_{M_{\setminus L}}(X_K,X_W \mid \doit(X_J)) = (g_K,\id_{\CXW})_*(P)$$
  since by Lemma~\ref{lem:marginalization_preserves_unique_solvability}, $g_K$ is the (unique) solution function of $M_{\setminus L}$.
  This means that both push-forwards are identical.

  Let $T \subseteq K \cup W$. Then $(M_{\setminus L})_{\doit(T)} = (M_{\doit(T)})_{\setminus L}$ by
  Proposition~\ref{prp:interventions_marginalizations_commute}. 
  The observational equivalence of $M_{\doit(T)}$ and $(M_{\doit(T)})_{\setminus L}$ w.r.t.\ $(K \cup W) \setminus T$ hence implies the observational equivalence of $M_{\doit(T)}$ and $(M_{\setminus L})_{\doit(T)}$ w.r.t.\ $(K \cup W) \setminus T$. 
  Since this holds for all $T \subseteq K \cup W$, $M$ and $M_{\setminus L}$ are interventionally equivalent w.r.t.\ $K \cup W$.
%  Finally, $(M_{\setminus L})^\twin = (M^\twin)_{\setminus (L \cup L')}$ by Proposition~\ref{prp:twinning_marginalization_commute}. Since $M^\twin$ and its marginalization
%  $(M^\twin)_{\setminus (L \cup L')}$ are interventionally equivalent w.r.t.\ $M \cup M' \cup W$, 
%  $M$ and $M_{\setminus L}$ are counterfactually equivalent w.r.t.\ $M \cup W$.
\end{proof}
So the marginalization operation indeed effectively hides the details of a subsystem, while preserving the causal semantics on the remaining part.
